% Post hoc checks -- regenerate with:
%   2-Dissertation-Draft-Works/wallet-reputation-experiments/scripts/run_supplementary_checks.py
% Generated by export_latex_results.py -- do not edit by hand unless re-export fails
\begin{table}[htbp]
\centering
\caption[스펜더 1{,}335개 가운데, 동결일 전에 다른 주소에서 송금을 받은 적이 있는 977개로 좁힌 스펜더 홀드아웃.]{스펜더 1{,}335개 가운데, 동결일 전에 다른 주소에서 송금을 받은 적이 있는 977개로 좁힌 스펜더 홀드아웃이에요. 나머지 358개에는 들어오는 송금 화살표가 없고, 스펜더 318개는 AWP에서 점수가 0이에요. 이렇게 좁히는 것은 결과를 본 뒤에 해 보았어요. 이 묶음에서 엔도스랭크와 AWP 사이의 차이는 표~\ref{tab:endorserank-awp-holdout}에 있어요. 신뢰구간은 지갑을 다시 뽑아 구한 95\% 백분위 부트스트랩 구간이에요(다시 뽑기 400번, 시드 42).}
\label{tab:holdout-receiving}
\fitwidth{%
\begin{tabular}{lcc}
\toprule
$t_1$의 점수 & 새 허락 쌍($\tau$ [95\% 신뢰구간]) & 새 송금자($\tau$ [95\% 신뢰구간]) \\
\midrule
EndorseRank & 0.340 [0.299, 0.386] & 0.358 [0.316, 0.396] \\
AWP & 0.298 [0.252, 0.343] & 0.388 [0.348, 0.430] \\
\midrule
$t_1$의 받은 허락 수 & 0.679 [0.616, 0.736] & 0.570 [0.505, 0.635] \\
$t_1$의 송금 입차수 & 0.569 [0.510, 0.631] & 0.701 [0.651, 0.750] \\
\midrule
AWP 빼기 t1 입차수(새 송금자) & \multicolumn{2}{c}{$\Delta\tau=-0.314$ [-0.353, -0.276]} \\
\bottomrule
\end{tabular}
}
\end{table}
